\section*{AI-use statement}

This paper was drafted by a crew of AI research-assistant agents
(Free-Research-Crew, protocol v0.1) under the direction of the human author, Bryce Wang.
The human approved the research question, the design and the results at three gates.
The submission gate is pending at the time of writing. The human author takes
responsibility for the content.

\begin{itemize}
  \item \textbf{Roles and models.} Every role subagent ran on Claude Haiku 5.5
        (\texttt{claude-haiku-5-5}): Scout, Data RA, Designer, Econometrician, Figures RA,
        Writer and Replicator, and the Reviewer. The PI session ran on Claude Opus 5.5.
  \item \textbf{Review.} The Reviewer agent ran 15 rounds: question (2 rounds, block then
        revise), design (3: revise, revise, pass), code (3: pass, revise, pass), results
        (3: revise, revise, pass) and paper (4: revise, revise, revise, revise). Every round
        was fresh-context. The Reviewer ran on the same model family as the authoring agents,
        so the independence level is fresh-context, not cross-family. The paper stage has been
        reviewed four times, and none of the four rounds passed. The reviews and the
        Reviewer's memory are in \texttt{reviews/}.
  \item \textbf{Findings accepted over the Reviewer: one.} design-r1 F4 (pre-gate
        computations in \texttt{design-mde}), kept by the PI's ruling with a disclosure,
        withdrawn by the Reviewer at design-r2 and design-r3, and shown to the human at the
        design gate.
  \item \textbf{Open findings at the question gate.} The human approved the question gate at
        04:22:30Z, with question-r2 F4, F7 and F8 deferred to the design stage, and question-r2
        F9 answered by the PI. The question stage was not re-reviewed after that, so those four
        records still read open. The design stage addressed their substance: the MDEs and the
        non-rejection wording (F4), the data checks (F7), a data check (F8), and the one-sided
        placebo cutoffs (F9). F8 is only partly met. The CSV's non-outcome moments match
        CCFT's printed values, but the CSV is not the \texttt{.dta} file CCFT used, and the
        moments of \texttt{vote} were not compared before the H1 run.
  \item \textbf{Clean reruns.} Four clean reruns, each from an environment built from
        \texttt{code/}\allowbreak\texttt{requirements.txt} alone. The first reproduced 7 of 15
        runs; its failures came from the rerun tool (it omitted \texttt{runs/} and counted a NaN
        in both runs as a mismatch), and the PI session fixed the tool (commit
        \texttt{0a72b72}). The second reproduced 15 of 15 (commit \texttt{c1c9549}), but its
        \texttt{replication.json} recorded the tool's interpreter at the top level. The PI
        session then changed the tool (commit \texttt{87890ae}) to record each run's interpreter
        and label the top-level value \texttt{crew\_python}. The third reproduced 15 of 15 with
        Python 3.12.11 recorded for every run. The fourth reproduced 15 of 15, with
        \texttt{crew.py} as of its last change (commit \texttt{87890ae}) and the workspace at
        commit \texttt{a6d7c45}. It was rerun because two claim texts were reworded (values
        unchanged) and two tables were regenerated, which made the third record older than its
        inputs. The NaN rule covered two metrics, \texttt{T\_Q\_asy} and \texttt{p\_Q\_asy}, which
        are NaN in both runs and so are reproduced without being compared. The fourth run's
        \texttt{replication.json} is the record the submission gate reads. No analysis code,
        data or tolerance changed. \texttt{REPLICATION.md} describes all four runs.
  \item \textbf{Briefs.} The role briefs are archived in \texttt{briefs/}.
  \item \textbf{Evidence.} Every reported number is generated from a logged run through
        the macro file \texttt{numbers.tex}. Every citation was checked against its source:
        four from the full text, and Eggers et al. (2015) from its abstract only. The full
        text of Cattaneo, Frandsen and Titiunik (2015) was read from a third-party mirror on
        GitHub (nbviewer). Its printed title, authors and DOI match the publisher record,
        and \texttt{sources/sources.json} records the DOI. The full workspace is
        \texttt{gallery/ss-rd-senate} in the Auto-Research-Skills repository; the public URL
        is to be added at release.
\end{itemize}
